\chapter*{\centering Abstract}
\addcontentsline{toc}{chapter}{Abstract}

Metabolic dysfunction-associated steatotic liver disease (MASLD) is managed
mainly through diet and weight loss, yet routine care offers neither a
personal estimate of liver fat nor recipes chosen for the individual.
HepatoTwin is a patient-specific digital twin that estimates liver fat from
routine check-up data, turns that estimate into evidence-linked nutritional
priorities, and ranks real recipes against them. A calibrated model trained on
15,300 NHANES participants with FibroScan measurements predicts the chance of
steatosis (AUC 0.881 for CAP $\geq$ 248~dB/m and 0.879 for CAP $\geq$
280~dB/m, calibration error $\leq$ 0.022), with conformal intervals that
covered 90.9\% of held-out cases. On NHANES III, it reached an AUC of 0.771
against ultrasound and identified a group with 1.58 times the all-cause and
2.15 times the cardiovascular mortality. An evidence engine links each
priority to a graded evidence table and a causal knowledge graph. A
health-gated, multi-criteria ranking of RecipeDB recipes placed recipes
meeting all three criteria in 82.3\% of top-ten lists, against 27.2\% for
weighted scoring without the gate, and produced complete seven-day plans for
every test patient. At each monthly check-in, the twin compares observed
weight with its own forecast and adjusts the plan. In a simulation calibrated
to published trials, its forecasts erred by 1.9~kg at three months, against
2.9~kg for a standard textbook rule. Patient summaries were fully faithful to
the model output. A web dashboard serves patients and clinicians, with plan
approval for safety concerns and a three-dimensional view of projected liver
change.

\vspace{0.4in}
\noindent\textbf{Keywords}: digital twin, MASLD, liver fat, NHANES,
personalised nutrition, recipe recommendation, explainable AI
\par
